\chapter{Spectral Clustering and Partitioning}

\chapterline

% --- local shorthands for this chapter (operator names only) ---
\providecommand{\cut}{\operatorname{cut}}
\providecommand{\vol}{\operatorname{vol}}
\providecommand{\diag}{\operatorname{diag}}
\providecommand{\RatioCut}{\operatorname{RatioCut}}
\providecommand{\NCut}{\operatorname{NCut}}

The previous two chapters developed the algebra of the graph Laplacian: the matrix
$\Lap = \mathbf{D} - \Adjm$, its non-negative eigenvalues $0 = \lambda_1 \le \lambda_2 \le \dots \le \lambda_n$,
and the quadratic form
\[
  x^\top \Lap x \;=\; \sum_{(i,j)\in E} w_{ij}(x_i - x_j)^2
\]
which measures how steeply a scalar assignment $x$ varies across connected edges. In this chapter, we put
that algebra to work on one of the central challenges in network analysis: \emph{graph partitioning}.
Given a graph, how can we partition its vertices into clusters that are densely connected on the inside
yet sparsely linked to the outside? Such clusters correspond to natural communities in a social network,
coherent objects in computer vision, or functional modules in biological networks.

Finding an optimal balanced cut in a graph is a notoriously difficult, NP-hard combinatorial puzzle.
Yet, an extraordinary geometric shortcut exists. If we relax the discrete requirement that vertices
choose between rigid binary labels and instead allow them to sit along continuous coordinates, the optimal
placement is governed by the bottom eigenvectors of the Laplacian. By computing these eigenvectors in
polynomial time, we can embed the graph into low-dimensional Euclidean space, where the natural clusters
pull apart into easily separable geometric point clouds.

\begin{intuitionbox}
Imagine pinning beads onto a taut wire, with elastic springs connecting any two beads that share an edge.
Each spring exerts a restoring force pulling its endpoints together, with total potential energy
$\sum_{(i,j)\in E}(x_i - x_j)^2$. If we simply release the beads, every bead rushes to a single point:
an energy of zero that collapses all structure.

To extract meaningful geometry, we impose two physical constraints: the center of mass must remain fixed
at zero ($\sum_i x_i = 0$), and the total spread must remain fixed ($\sum_i x_i^2 = n$). Under these
rules, the beads cannot collapse. Instead, vertices that are tightly connected by many internal springs
huddle close together, while the few sparse bridges between clusters are forced to stretch across the gap.
The equilibrium positions that emerge are precisely the coordinates of the \emph{Fiedler vector} $v_2$,
the eigenvector of $\Lap$ corresponding to its second-smallest eigenvalue $\lambda_2$. Slicing the wire
at a suitable threshold reveals the bottleneck cut.
\end{intuitionbox}

We will develop this geometric viewpoint in three systematic stages:
\begin{enumerate}
  \item \textbf{Spectral Bisection}: Using the single Fiedler vector $v_2$ to project vertices onto a line,
        followed by a sweep threshold to find a clean balanced two-way cut.
  \item \textbf{Normalized Cuts and Random Walks}: Normalizing by degree matrices so that cluster size
        reflects edge volume rather than bare vertex count, linking the cut directly to random-walk escape
        probabilities and Cheeger's inequality.
  \item \textbf{$k$-Way Spectral Clustering}: Projecting vertices into $\R^k$ using the bottom $k$
        eigenvectors simultaneously, where communities become well-separated point clouds clustered by
        standard $k$-means.
\end{enumerate}
Throughout, we trace every step by hand on a concrete running network so that the linear algebra and the
geometry remain visible at every turn.

\section{The cut we wish we could compute}
\label{ch33:the-cut}

Fix an undirected graph $\Gph=(V,E)$ with nonnegative edge weights $w_{ij}\ge 0$. For disjoint vertex sets
$A,B\subseteq V$ define the cut weight
\[
  \cut(A,B) \;=\; \sum_{i\in A,\,j\in B} w_{ij},
\]
and for a single set $S$ write $\cut(S)=\cut(S,\overline S)$ where $\overline S = V\setminus S$. The
\emph{minimum cut} -- smallest $\cut(S)$ over all nonempty proper $S$ -- is computable in polynomial time
(Chapter~11), but it is useless for clustering: the cheapest cut usually peels off a single low-degree
vertex. We do not want \emph{a} cut; we want a \emph{balanced} cut, one that genuinely splits the graph.

Two classical balance-aware objectives. With $\vol(S)=\sum_{i\in S}d_i$ the volume (sum of weighted
degrees) and $\abs{S}$ the cardinality:
\[
  \RatioCut(A,B) \;=\; \cut(A,B)\!\left(\frac{1}{\abs{A}}+\frac{1}{\abs{B}}\right),
  \qquad
  \NCut(A,B) \;=\; \cut(A,B)\!\left(\frac{1}{\vol(A)}+\frac{1}{\vol(B)}\right).
\]

\begin{focusbox}
Both objectives multiply the raw cut by a penalty that \emph{blows up when one side is small}. Peeling off
one vertex makes $\abs A=1$ (or $\vol(A)=d$ small), so the factor $1/\abs A$ (or $1/\vol(A)$) is large and
the objective is bad even though $\cut$ is tiny. The minimizer is forced to keep both sides fat. RatioCut
balances by \emph{count}; normalized cut balances by \emph{volume}. They differ only in which notion of
``size'' they protect.
\end{focusbox}

\begin{theorem}[Hardness]
\label{ch33:hardness}
Minimizing $\RatioCut$ over all bipartitions of $V$ is NP-hard, as is minimizing $\NCut$.
\end{theorem}

Because exact combinatorial optimization is NP-hard, we turn to continuous relaxation. Rather than
forcing each coordinate to take one of two discrete values, we allow the coordinates to vary freely
along the real line. The combinatorial objective then transforms into a continuous Rayleigh quotient
involving the graph Laplacian $\Lap$. Solving this relaxation amounts to computing an eigenvector, which
can be done in polynomial time. Once we have the continuous solution, we round the real-valued coordinates
back to a discrete cut.

\subsection*{Our running graph}

\begin{figure}[H]
\centering
\begin{tikzpicture}[
   v/.style={circle,draw=nodeblue,fill=nodeblue!12,thick,minimum size=8mm,inner sep=0pt,font=\small\bfseries},
   vr/.style={circle,draw=targetgreen,fill=targetgreen!12,thick,minimum size=8mm,inner sep=0pt,font=\small\bfseries},
   intra/.style={draw=edgegray,thick},
   bridge/.style={draw=trapred,very thick}]
  % Cluster background hulls
  \draw[rounded corners=10pt,fill=sourceblue!6,draw=sourceblue!30,thick] (-0.5,-0.1) rectangle (1.7,2.5);
  \node[font=\footnotesize\bfseries,text=sourceblue] at (0.6,2.8) {Cluster $L = \{1,2,3\}$};

  \draw[rounded corners=10pt,fill=targetgreen!6,draw=targetgreen!30,thick] (3.7,-0.1) rectangle (5.9,2.5);
  \node[font=\footnotesize\bfseries,text=targetgreen] at (4.8,2.8) {Cluster $R = \{4,5,6\}$};

  % left cluster (triangle): 1,2,3
  \node[v] (1) at (0,1.2)   {1};
  \node[v] (2) at (1.1,2.0) {2};
  \node[v] (3) at (1.1,0.4) {3};
  % right cluster: 4,5,6
  \node[vr] (4) at (4.3,1.2) {4};
  \node[vr] (5) at (5.4,2.0) {5};
  \node[vr] (6) at (5.4,0.4) {6};

  % intra-cluster edges
  \path[intra] (1) edge (2) (1) edge (3) (2) edge (3);
  \path[intra] (4) edge (5) (4) edge (6) (5) edge (6);

  % the single bridge
  \path[bridge] (3) edge node[above=2pt,font=\scriptsize\bfseries,text=trapred] {bridge $(3,4)$} (4);

  % cut partition boundary
  \draw[trapred!80,thick,dashed] (2.7,-0.3) -- (2.7,2.5);
  \node[font=\scriptsize\bfseries,text=trapred,fill=white,inner sep=1.5pt,rounded corners=2pt] at (2.7,-0.5) {bipartition cut ($\cut = 1$)};
\end{tikzpicture}
\caption{The running graph $\Gph_0$: two dense triangles joined by a single sparse bridge $(3,4)$. All edge
weights are $1$. The target cut is $L=\{1,2,3\} \mid R=\{4,5,6\}$ with cut weight $1$, achieving optimal
balance and conductance.}
\label{ch33:fig-running}
\end{figure}

$\Gph_0$ has $n=6$ vertices and $7$ edges. The degrees are $d_1=d_2=d_5=d_6=2$ and $d_3=d_4=3$ (the
bridge endpoints carry one extra edge). The intended answer is obvious to the eye -- the two triangles --
and that is exactly why it is a good test: a correct algorithm must agree with the eye, and we can check
every number by hand.

\section{Bisection by the Fiedler vector}
\label{ch33:fiedler}

Take the RatioCut objective with the soft relaxation. Encode a bipartition $(A,B)$ by the vector
\[
  f_i \;=\;
  \begin{cases}
    +\sqrt{\abs{B}/\abs{A}}, & i\in A,\\[2pt]
    -\sqrt{\abs{A}/\abs{B}}, & i\in B.
  \end{cases}
\]
This unusual scaling is chosen so that two clean identities hold. A short computation (do it in the margin)
gives
\[
  \sum_i f_i = \abs A\sqrt{\tfrac{\abs B}{\abs A}} - \abs B\sqrt{\tfrac{\abs A}{\abs B}} = 0,
  \qquad
  \norm{f}^2 = \abs A\cdot\tfrac{\abs B}{\abs A} + \abs B\cdot\tfrac{\abs A}{\abs B} = n,
\]
so $f\perp \mathbf 1$ and $\norm f^2 = n$. The point of the scaling is the next identity.

\begin{lemma}[RatioCut as a quadratic form]
\label{ch33:ratiocut-quad}
With $f$ defined above, $\; f^\top \Lap f = n\cdot\RatioCut(A,B).$
\end{lemma}

\begin{proof}
Using $x^\top\Lap x=\sum_{(i,j)\in E}w_{ij}(x_i-x_j)^2$, only edges crossing the cut contribute, and a
crossing edge has $(f_i-f_j)^2 = \big(\sqrt{\abs B/\abs A}+\sqrt{\abs A/\abs B}\big)^2$. Expand:
\[
  \Big(\sqrt{\tfrac{\abs B}{\abs A}}+\sqrt{\tfrac{\abs A}{\abs B}}\Big)^2
  = \tfrac{\abs B}{\abs A} + 2 + \tfrac{\abs A}{\abs B}
  = \tfrac{\abs A+\abs B}{\abs A} + \tfrac{\abs A+\abs B}{\abs B}
  = n\Big(\tfrac1{\abs A}+\tfrac1{\abs B}\Big).
\]
Summing over the $\cut(A,B)$ crossing edges gives $f^\top\Lap f = \cut(A,B)\cdot n(\tfrac1{\abs A}+\tfrac1{\abs B})
= n\cdot\RatioCut(A,B)$.
\end{proof}

So minimizing RatioCut is \emph{exactly}
\[
  \min_{A,B}\; f^\top \Lap f
  \quad\text{subject to}\quad
  f \text{ of the two-valued form above},\ \ f\perp\mathbf 1,\ \ \norm f^2=n.
\]
The constraint forcing $f$ to take only two discrete values is the sole source of combinatorial hardness.
To overcome it, we relax the problem: we drop the two-valued requirement and allow $f\in\R^n$ to be an
arbitrary real vector, retaining only the centering constraint $f\perp\mathbf 1$ and the scale constraint
$\norm f^2 = n$:
\[
  \min_{f\in\R^n}\ f^\top\Lap f
  \quad\text{subject to}\quad f\perp\mathbf 1,\ \norm f^2 = n.
\]

\begin{whybox}
\textbf{The Rayleigh quotient and the Courant--Fischer theorem.}
For any symmetric matrix $\mathbf{M}$ with eigenvalues $\lambda_1 \le \lambda_2 \le \dots \le \lambda_n$
and orthonormal eigenvectors $u_1, \dots, u_n$, the \emph{Rayleigh quotient}
\[
  R_{\mathbf{M}}(x) \;=\; \frac{x^\top \mathbf{M} x}{x^\top x}
\]
measures the quadratic energy per unit norm. Expanding an arbitrary vector $x$ in the eigenvector basis as
$x = \sum_{i=1}^n c_i u_i$ gives $\norm{x}^2 = \sum_i c_i^2$ and $x^\top \mathbf{M} x = \sum_i \lambda_i c_i^2$.
The Rayleigh quotient is therefore an exact weighted average of the eigenvalues:
\[
  R_{\mathbf{M}}(x) \;=\; \frac{\sum_{i=1}^n \lambda_i c_i^2}{\sum_{i=1}^n c_i^2}.
\]
Because the weights $c_i^2$ are non-negative, this weighted average can never be smaller than the minimum
eigenvalue $\lambda_1$, nor larger than the maximum $\lambda_n$. The global minimum $\lambda_1$ is attained
by choosing $x = u_1$.

What happens if we constrain $x$ to be orthogonal to the first eigenvector $u_1$? That condition forces
$c_1 = \langle x, u_1 \rangle = 0$. In the weighted average, the term $\lambda_1 c_1^2$ vanishes, leaving only
eigenvalues $\lambda_2, \lambda_3, \dots, \lambda_n$. The minimum of this restricted average is unmistakably
$\lambda_2$, achieved when $x = u_2$.

More generally, the \emph{Courant--Fischer min-max theorem} states that for each index $k$,
\[
  \lambda_k \;=\; \min_{\dim S = k}\; \max_{0\ne x\in S}\; R_{\mathbf{M}}(x)
  \;=\; \min_{x \perp u_1, \dots, u_{k-1},\, x\ne 0} R_{\mathbf{M}}(x).
\]
Applying this to the graph Laplacian $\Lap$:
\begin{itemize}
  \item The smallest eigenvalue is always $\lambda_1 = 0$, with eigenvector $u_1 = \frac{1}{\sqrt{n}}\mathbf{1}$
        (since $\Lap \mathbf{1} = 0$).
  \item The constraint $f \perp \mathbf{1}$ in our relaxed bisection is precisely $f \perp u_1$.
\end{itemize}
By Courant--Fischer, minimizing $f^\top \Lap f$ over all vectors orthogonal to $\mathbf{1}$ with $\norm{f}^2 = n$
yields exactly $n \lambda_2$, and the optimal direction is the second eigenvector $v_2$. This eigenvector is
the \emph{Fiedler vector}, named after Miroslav Fiedler who in 1973 identified $\lambda_2$ as the
\emph{algebraic connectivity} of the graph.
\end{whybox}

We have proved:

\begin{theorem}[Fiedler relaxation]
\label{ch33:fiedler-relax}
The continuous relaxation of RatioCut bisection is solved by the Fiedler vector $v_2$, with optimal value
$\lambda_2 = n^{-1}\min f^\top\Lap f$. Consequently $\RatioCut^\star \ge \lambda_2$: \emph{the algebraic
connectivity is a lower bound on the best ratio cut.}
\end{theorem}

Now we must \emph{round} the real vector $v_2$ back to a partition. The natural move: pick a threshold $t$
and set $A=\set{i: (v_2)_i > t}$, $B$ the rest. Two standard thresholds:
\begin{itemize}
  \item \textbf{Sign threshold} ($t=0$): split by the sign of $(v_2)_i$. Simple; works when the gap straddles
        zero.
  \item \textbf{Median / sweep threshold}: sort the entries of $v_2$, and among the $n-1$ ``sweep cuts''
        (prefixes of the sorted order) pick the one with the smallest RatioCut. This is the \emph{sweep},
        and it is what good implementations use.
\end{itemize}

\begin{algorithm}[H]
\caption{\textsc{SpectralBisection}$(\Gph)$ -- Fiedler partition with sweep rounding}
\label{ch33:alg-bisection}
\DontPrintSemicolon
\KwIn{Connected weighted graph $\Gph=(V,E,w)$, $n=\abs V$.}
\KwOut{A bipartition $(A,B)$ of $V$.}
Form the Laplacian $\Lap = \mathbf D - \Adjm$.\;
Compute the eigenvector $v_2$ for the second-smallest eigenvalue $\lambda_2$.\;
Sort vertices so that $(v_2)_{\pi(1)} \le (v_2)_{\pi(2)} \le \dots \le (v_2)_{\pi(n)}$.\;
$best \gets +\infty$;\quad $A^\star\gets\emptyset$.\;
\For{$k \gets 1$ \KwTo $n-1$}{
   $A \gets \set{\pi(1),\dots,\pi(k)}$,\quad $B \gets V\setminus A$.\;
   $q \gets \RatioCut(A,B)$ \tcp*{or $\NCut$ for the normalized variant}
   \If{$q < best$}{ $best \gets q$;\quad $A^\star \gets A$. }
}
\Return $(A^\star,\, V\setminus A^\star)$.\;
\end{algorithm}

\begin{trapbox}
The eigenvalue solver returns $v_2$ up to sign and scale. The sign is meaningless -- $-v_2$ is just as
valid an eigenvector -- so \emph{never} attach semantic meaning to ``positive = cluster A.'' Always round
through a sweep or a sign split and read off the partition, never the raw signs as labels across two runs.
A second trap: if the graph is \emph{disconnected}, $\lambda_2=0$ and $v_2$ is not unique; the relaxation
silently returns a connected component, which may or may not be what you wanted. Always check connectivity
first ($\lambda_2 > 0 \iff \Gph$ connected, the spectral characterization from the Laplacian chapter).
\end{trapbox}

\subsection*{Tracing $\Gph_0$ by hand}
\label{ch33:trace-bisection}

The Laplacian of the running graph (order $1,2,3,4,5,6$):
\[
  \Lap =
  \begin{pmatrix}
    2 & -1 & -1 & 0 & 0 & 0\\
    -1 & 2 & -1 & 0 & 0 & 0\\
    -1 & -1 & 3 & -1 & 0 & 0\\
    0 & 0 & -1 & 3 & -1 & -1\\
    0 & 0 & 0 & -1 & 2 & -1\\
    0 & 0 & 0 & -1 & -1 & 2
  \end{pmatrix}.
\]
The graph has a mirror symmetry swapping $L=\set{1,2,3}$ with $R=\set{4,5,6}$ (send $1\!\leftrightarrow\!6$,
$2\!\leftrightarrow\!5$, $3\!\leftrightarrow\!4$). Eigenvectors are either symmetric or antisymmetric under
this swap. The Fiedler vector is the antisymmetric one, constant on each triangle except at the bridge
endpoints. Solving the $6\times6$ system (or just trusting the symmetry and a $2$-parameter ansatz
$f = (a,a,b,-b,-a,-a)$) yields, after normalizing,
\[
  v_2 \;\approx\; (\,0.46,\ 0.46,\ 0.26,\ -0.26,\ -0.46,\ -0.46\,)^\top,
  \qquad \lambda_2 \approx 0.44.
\]

\begin{figure}[H]
\centering
\begin{tikzpicture}[scale=1.1,
   vL/.style={circle,draw=sourceblue,fill=sourceblue!15,thick,inner sep=0pt,minimum size=6.5mm,font=\scriptsize\bfseries},
   vR/.style={circle,draw=targetgreen,fill=targetgreen!15,thick,inner sep=0pt,minimum size=6.5mm,font=\scriptsize\bfseries}]

  % Background shaded bands for clusters
  \fill[targetgreen!6,rounded corners=4pt] (-0.8,-1.2) rectangle (2.85,1.2);
  \fill[sourceblue!6,rounded corners=4pt] (3.05,-1.2) rectangle (6.7,1.2);

  % Axis line
  \draw[->,thick,edgegray] (-0.8,0) -- (6.8,0) node[right,font=\small]{$(v_2)_i$};

  % Coordinates mapping: -0.46 -> 0.4, -0.26 -> 1.7, 0.0 -> 2.95, +0.26 -> 4.2, +0.46 -> 5.5
  % Tick marks
  \foreach \x/\lbl in {0.4/$-0.46$, 1.7/$-0.26$, 2.95/$0$, 4.2/$+0.26$, 5.5/$+0.46$} {
    \draw[thick,edgegray] (\x,0.1) -- (\x,-0.1);
    \node[below=4pt,font=\scriptsize,text=edgegray] at (\x,0) {\lbl};
  }

  % Vertices plotted at their coordinates
  \node[vR] at (0.4,0.65) {5};
  \node[vR] at (0.4,-0.65) {6};
  \node[vR] at (1.7,0.5) {4};

  \node[vL] at (4.2,0.5) {3};
  \node[vL] at (5.5,0.65) {1};
  \node[vL] at (5.5,-0.65) {2};

  % Threshold cut line at 0
  \draw[trapred,very thick,dashed] (2.95,-1.15) -- (2.95,1.15)
        node[above=2pt,font=\scriptsize\bfseries,text=trapred]{threshold $t=0$};

  % Cluster labels and sweep indications
  \node[font=\footnotesize\bfseries,text=targetgreen] at (1.0,1.6) {Cluster $R$ (negative Fiedler coords)};
  \node[font=\footnotesize\bfseries,text=sourceblue] at (4.9,1.6) {Cluster $L$ (positive Fiedler coords)};

  % Sweep arrows showing candidate prefix order
  \draw[->,thick,quotientorange] (0.0,-1.4) -- (5.9,-1.4)
        node[midway,below=2pt,font=\scriptsize\bfseries]{Sweep order: $5 \prec 6 \prec 4 \;\mid\; 3 \prec 1 \prec 2$};
\end{tikzpicture}
\caption{Vertices of $\Gph_0$ embedded along the 1D line by their Fiedler coordinates $(v_2)_i$. The dense
triangles collapse to tight groups on opposite sides of the origin, leaving a wide gap around $0$.
Thresholding at $t=0$ cleanly isolates cluster $R=\{4,5,6\}$ from cluster $L=\{1,2,3\}$.}
\label{ch33:fig-line}
\end{figure}

Run the sweep. Sorted order is $5,6,4,\,|\,3,1,2$ (ties broken arbitrarily inside a triangle). The sweep
cuts and their RatioCuts:

\begin{center}\small
\begin{tabular}{c l c c l}
\toprule
$k$ & $A$ (smallest $k$) & $\cut$ & $\frac1{\abs A}+\frac1{\abs B}$ & RatioCut\\
\midrule
1 & $\set5$           & 2 & $1 + \tfrac15$ & $2.40$\\
2 & $\set{5,6}$       & 2 & $\tfrac12+\tfrac14$ & $1.50$\\
\textbf 3 & $\set{4,5,6}$ & \textbf 1 & $\tfrac13+\tfrac13$ & $\mathbf{0.667}$\\
4 & $\set{3,4,5,6}$   & 2 & $\tfrac14+\tfrac12$ & $1.50$\\
5 & $\set{1,3,4,5,6}$ & 2 & $\tfrac15+1$ & $2.40$\\
\bottomrule
\end{tabular}
\end{center}

The minimum is at $k=3$: $A=\set{4,5,6}$, RatioCut $=0.667$, exactly the two triangles. The lower bound from
Theorem~\ref{ch33:fiedler-relax} was $\lambda_2\approx0.44 \le 0.667$ -- the relaxation undershoots the true
optimum, as it must, and the gap measures how much rounding cost we paid.

\begin{labbox}
Redo this on the \emph{path} $P_4$ (vertices $1\!-\!2\!-\!3\!-\!4$). Its Fiedler vector is
$v_2\propto(1.31,\,0.54,\,-0.54,\,-1.31)$ and $\lambda_2 = 2-\sqrt2\approx0.586$. The sign split gives
$\set{1,2}\mid\set{3,4}$ -- cut the middle edge, the only sensible bisection of a path. Confirm the sweep
agrees, then try $P_5$ and watch the middle vertex land near zero, genuinely ambiguous.
\end{labbox}

\section{Why the relaxation is good: gaps and Cheeger}
\label{ch33:cheeger}

A relaxation is only worth trusting if the rounded answer is provably close to optimal. For balanced cuts
the guarantee is \emph{Cheeger's inequality}, the central theorem of spectral partitioning. Define the
\emph{conductance} (a normalized, balance-aware cut quality) of a set $S$ and of the graph:
\[
  \phi(S) \;=\; \frac{\cut(S)}{\min\set{\vol(S),\,\vol(\overline S)}},
  \qquad
  \phi(\Gph) \;=\; \min_{S\ne\emptyset,\,S\ne V} \phi(S).
\]
Small conductance means a sparse, balanced cut -- precisely a good cluster boundary. Let $\lambda_2^{\mathrm{nrm}}$
be the second-smallest eigenvalue of the \emph{normalized} Laplacian $\Lap_{\mathrm{sym}} =
\mathbf D^{-1/2}\Lap\,\mathbf D^{-1/2}$ (introduced next section).

\begin{theorem}[Cheeger's inequality]
\label{ch33:cheeger-thm}
For any connected graph,
\[
  \frac{\lambda_2^{\mathrm{nrm}}}{2} \;\le\; \phi(\Gph) \;\le\; \sqrt{2\,\lambda_2^{\mathrm{nrm}}}.
\]
\end{theorem}

\begin{whybox}
The left inequality (the easy direction) says the eigenvalue lower-bounds the best conductance: no cut can
be sparser than $\lambda_2^{\mathrm{nrm}}/2$, mirroring Theorem~\ref{ch33:fiedler-relax}. The right
inequality is the algorithmic gold: it is \emph{constructive}. Its proof takes the eigenvector $v_2$, runs
exactly the \emph{sweep} of Algorithm~\ref{ch33:alg-bisection}, and proves that some sweep cut achieves
conductance $\le\sqrt{2\lambda_2^{\mathrm{nrm}}}$. So the simple ``sort by eigenvector, try every prefix''
rounding is not a heuristic -- it provably finds a cut within a quadratic factor of optimal.
\end{whybox}

\begin{proof}[Proof sketch of the hard direction]
Let $g=\mathbf D^{1/2}v_2$, so the Rayleigh quotient
$R(g)=\frac{g^\top\Lap g}{\sum_i d_i g_i^2}=\lambda_2^{\mathrm{nrm}}$. Order the vertices by $g$ and
consider the $n-1$ sweep sets $S_t = \set{i : g_i > t}$. The key combinatorial lemma: if \emph{every} sweep
cut had conductance $>\phi$, then a probabilistic argument over a random threshold $t$ (drawn with density
proportional to $2\abs t$) forces $R(g) > \phi^2/2$, contradicting $R(g)=\lambda_2^{\mathrm{nrm}}$ once we
set $\phi=\sqrt{2\lambda_2^{\mathrm{nrm}}}$. Concretely, write the numerator and denominator of $R(g)$ as
expectations over the random threshold; Cauchy--Schwarz relates the crossing-edge probability to the
conductance of the realized sweep cut, and the bound on the minimum over sweep cuts pops out. The full
argument is in the spectral graph theory chapter; the takeaway is that \emph{the sweep over the eigenvector
is the rounding the theorem analyzes.}
\end{proof}

\begin{connectionbox}
Cheeger's inequality is the graph incarnation of a Riemannian-geometry theorem (Cheeger 1970, Buser, Alon--
Milman in the discrete setting): the smallest nonzero eigenvalue of the Laplace operator controls the
narrowest ``waist'' of the space. On graphs it ties three worlds together -- \textbf{spectra} (linear
algebra book: eigenvalues of symmetric matrices), \textbf{cuts} (this book: max-flow / min-cut, expansion),
and \textbf{random walks} (probability book: $\lambda_2^{\mathrm{nrm}}$ is the spectral gap that governs
mixing time, Chapters~35--36). A graph is an \emph{expander} -- the subject of the next chapter -- exactly
when $\phi(\Gph)$ is bounded below by a constant, equivalently $\lambda_2^{\mathrm{nrm}}=\Om(1)$.
\end{connectionbox}

\section{Normalized cuts and the random-walk view}
\label{ch33:ncut}

RatioCut protects \emph{cardinality} balance; normalized cut protects \emph{volume} balance, which matches
real data where dense clusters should be measured by total connection mass, not vertex count. The relaxation
mirrors Section~\ref{ch33:fiedler} but with $\mathbf D$-weighted inner products.

Encode a partition by $h_i = \sqrt{\vol(B)/\vol(A)}$ on $A$ and $-\sqrt{\vol(A)/\vol(B)}$ on $B$. The same
algebra as Lemma~\ref{ch33:ratiocut-quad}, now with the weighting $\mathbf D$, gives
\[
  h^\top\Lap h = \vol(V)\cdot\NCut(A,B),\qquad h^\top\mathbf D h = \vol(V),\qquad (\mathbf D h)\perp\mathbf 1 .
\]
Substituting $h = \mathbf D^{-1/2}g$ turns the constrained problem into
\[
  \min_g\ \frac{g^\top \mathbf D^{-1/2}\Lap\,\mathbf D^{-1/2}\, g}{g^\top g}
  = \min_g\ \frac{g^\top \Lap_{\mathrm{sym}}\,g}{g^\top g}
  \quad\text{s.t. } g\perp \mathbf D^{1/2}\mathbf 1 .
\]
The minimizer is the second eigenvector of the \emph{symmetric normalized Laplacian}
$\Lap_{\mathrm{sym}}=\mathbf D^{-1/2}\Lap\,\mathbf D^{-1/2} = \mathbf I - \mathbf D^{-1/2}\Adjm\mathbf D^{-1/2}$,
and the cluster indicator is recovered as $h=\mathbf D^{-1/2}g$. There is a closely related
\emph{random-walk Laplacian} $\Lap_{\mathrm{rw}} = \mathbf D^{-1}\Lap = \mathbf I - \mathbf D^{-1}\Adjm$,
whose eigenvectors are exactly the $h$'s and whose matrix $\mathbf P=\mathbf D^{-1}\Adjm$ is the transition
matrix of the natural random walk.

\begin{intuitionbox}
Drop a random walker on the graph. It tends to get \emph{trapped} inside a dense cluster: many edges keep
it bouncing around inside, few edges let it escape. Normalized cut is, almost literally, the probability
that a stationary walker crosses the boundary in one step. Minimizing $\NCut$ finds the partition the walker
finds hardest to cross -- the bottleneck. This is why $\Lap_{\mathrm{rw}}$, the walk's generator, is the
right operator, and why its spectral gap is the walk's mixing rate.
\end{intuitionbox}

\begin{formalbox}
The three Laplacians share eigenvalues in a tidy way. If $\Lap_{\mathrm{sym}} g = \mu g$ then
$\Lap_{\mathrm{rw}}(\mathbf D^{-1/2}g) = \mu(\mathbf D^{-1/2}g)$ -- same eigenvalue $\mu\in[0,2]$, eigenvectors
related by $\mathbf D^{-1/2}$. Use $\Lap_{\mathrm{sym}}$ for computation (it is symmetric, so solvers are
stable and the eigenvectors are orthonormal) and $\Lap_{\mathrm{rw}}$ for interpretation (its eigenvectors
are the cluster indicators $h$, not the rescaled $g$). The unnormalized $\Lap$ corresponds to RatioCut;
the normalized pair corresponds to NCut. In practice the normalized version is preferred -- it is more
robust when degrees vary, a fact made precise by Cheeger's bound, which uses the \emph{normalized} gap.
\end{formalbox}

\subsection*{Tracing the normalized cut on $\Gph_0$}

For $\Gph_0$ the degrees are $\mathbf D=\diag(2,2,3,3,2,2)$. Because the graph is nearly regular, the
normalized Fiedler vector $g$ is close to the unnormalized one but tilted: the degree-$3$ bridge vertices
$3,4$ get pulled toward zero a little harder (division by $\sqrt3$). Numerically the second eigenvalue of
$\Lap_{\mathrm{sym}}$ is $\mu_2 \approx 0.20$, and the recovered indicator $h=\mathbf D^{-1/2}g$ again
changes sign exactly between $\set{1,2,3}$ and $\set{4,5,6}$. The NCut value of that split is
\[
  \NCut(L,R) = \cut(L,R)\Big(\tfrac1{\vol L}+\tfrac1{\vol R}\Big)
  = 1\cdot\Big(\tfrac17+\tfrac17\Big) = \tfrac27 \approx 0.286,
\]
since $\vol L=\vol R = 2+2+3 = 7$. Every other sweep cut scores worse (peeling a single degree-$2$ vertex
costs $2\cdot(\tfrac12 + \tfrac1{12})\approx1.17$), so the sweep again returns the two triangles. Cheeger's
lower bound $\mu_2/2\approx0.10 \le \phi(\Gph_0)$ is consistent with $\phi(\Gph_0)=\cut(L)/\vol(L)=1/7\approx0.143$.

\section{$k$-way spectral clustering}
\label{ch33:kway}

For $k>2$ clusters, recursively bisecting works but is myopic; the clean approach uses the \emph{bottom $k$
eigenvectors at once}. The idea generalizes the line embedding of Figure~\ref{ch33:fig-line} from one
dimension to $k$.

\begin{intuitionbox}
One eigenvector lays the graph on a line. Two eigenvectors lay it on a plane; $k-1$ of them lay it in
$\R^{k-1}$. Each vertex $i$ becomes a point $y_i\in\R^{k}$ whose coordinates are its entries in
$v_1,\dots,v_k$. Well-separated clusters in the graph become well-separated \emph{point clouds} in this
\emph{spectral embedding}, and we finish with ordinary geometric clustering -- $k$-means -- in that space.
The Laplacian has done the hard work of bending the graph into a shape where Euclidean distance means
``hard to reach by walking on the graph.''
\end{intuitionbox}

Why the bottom $k$? Consider the ideal case of $k$ \emph{disconnected} clusters. Then $\Lap$ is
block-diagonal, $\lambda_1=\dots=\lambda_k=0$, and the eigenspace of $0$ is spanned by the indicator vectors
$\mathbf 1_{C_1},\dots,\mathbf 1_{C_k}$ of the components. Any orthonormal basis of that eigenspace gives an
embedding in which \emph{all vertices of one cluster map to a single point} of $\R^k$ -- $k$ distinct points
total. $k$-means recovers the clusters exactly. When the clusters are merely well-separated rather than
disconnected, the eigenvalues $\lambda_2,\dots,\lambda_k$ lift slightly off $0$ and the points spread into
$k$ tight blobs; matrix perturbation theory (the Davis--Kahan $\sin\Theta$ theorem from the linear algebra
book) bounds how much the embedding moves, hence how clean $k$-means stays.

\begin{algorithm}[H]
\caption{\textsc{NormalizedSpectralClustering}$(\Gph,k)$ -- Ng--Jordan--Weiss}
\label{ch33:alg-kway}
\DontPrintSemicolon
\KwIn{Weighted graph $\Gph=(V,E,w)$, number of clusters $k$.}
\KwOut{A partition $V = C_1\cup\dots\cup C_k$.}
$\mathbf D \gets \diag(d_1,\dots,d_n)$;\quad $\Lap_{\mathrm{sym}} \gets \mathbf I - \mathbf D^{-1/2}\Adjm\,\mathbf D^{-1/2}$.\;
Compute the $k$ eigenvectors $u_1,\dots,u_k$ of $\Lap_{\mathrm{sym}}$ for its $k$ \emph{smallest} eigenvalues.\;
$U \gets [\,u_1\ \cdots\ u_k\,] \in \R^{n\times k}$ \tcp*{vertex $i$ is row $i$ of $U$}
\For{$i \gets 1$ \KwTo $n$}{
   $y_i \gets U_{i,\cdot} \,/\, \norm{U_{i,\cdot}}_2$ \tcp*{row-normalize onto the unit sphere}
}
$\set{C_1,\dots,C_k} \gets \textsc{$k$-Means}(\set{y_1,\dots,y_n},\,k)$.\;
\Return the partition assigning $i$ to its $k$-means cluster.\;
\end{algorithm}

\begin{trapbox}
The row-normalization on line 5 (project each embedded point to the unit sphere) is the Ng--Jordan--Weiss
trick and it matters: without it, high-degree and low-degree vertices land at different radii and $k$-means
splits by radius instead of by cluster. A second subtlety: you need the $k$ \emph{smallest} eigenvalues, so
use a solver tuned for the bottom of the spectrum (shift-and-invert Lanczos, or the Laplacian
solvers of Chapters~31--32), not a top-$k$ method. Finally, $k$-means is non-convex; run it with several
$k$-means++ seedings and keep the lowest-inertia result, or the embedding's good work can be wrecked by a
bad initialization.
\end{trapbox}

\subsection*{A worked $k=2$ embedding, and a $k=3$ sketch}

For $k=2$ on $\Gph_0$, $U=[\,u_1\ u_2\,]$ has first column $u_1\propto\mathbf D^{1/2}\mathbf 1$ (the constant
direction, same for everyone up to degree) and second column $u_2 = $ normalized Fiedler. After
row-normalization every vertex sits on the unit circle; the three $L$-vertices cluster near one point and
the three $R$-vertices near the antipodal point, because $u_2$ separates them while $u_1$ is shared. Two-means
immediately recovers $\set{1,2,3}$ and $\set{4,5,6}$. This is the same partition the sweep found, now obtained
geometrically -- reassuring, and the general mechanism for larger $k$.

To see $k=3$, glue a \emph{third} triangle $T_3=\set{7,8,9}$ to $\Gph_0$ with one bridge $(6,7)$, giving a
chain of three triangles. The bottom three eigenvalues are $0$ and two small positives; the embedding
$U\in\R^{9\times3}$ sends each triangle to a tight blob at a corner of a triangle in $\R^3$, and $3$-means
returns the three communities. The eigengap between $\lambda_3$ (small) and $\lambda_4$ (an order larger,
governed by \emph{within}-triangle structure) is the spectral signal that ``$k=3$ is the right number of
clusters'' -- the \emph{eigengap heuristic} for choosing $k$. Concretely the bottom of the spectrum is
$\lambda_1=0,\ \lambda_2\approx0.21,\ \lambda_3\approx0.70,\ \lambda_4=3$, so the jump from $\lambda_3$ to
$\lambda_4$ is the visible gap.

\begin{figure}[H]
\centering
\begin{tikzpicture}[
   v1/.style={circle,draw=sourceblue,fill=sourceblue!15,thick,minimum size=6.5mm,inner sep=0pt,font=\scriptsize\bfseries},
   v2/.style={circle,draw=targetgreen,fill=targetgreen!15,thick,minimum size=6.5mm,inner sep=0pt,font=\scriptsize\bfseries},
   v3/.style={circle,draw=quotientorange,fill=quotientorange!15,thick,minimum size=6.5mm,inner sep=0pt,font=\scriptsize\bfseries},
   intra/.style={draw=edgegray,thick},
   bridge/.style={draw=trapred,very thick}]

  % Background shaded panels for each community
  \draw[rounded corners=6pt,fill=sourceblue!6,draw=sourceblue!25,thick] (-0.5,-0.2) rectangle (1.5,2.1);
  \node[font=\scriptsize\bfseries,text=sourceblue] at (0.5,2.35) {$C_1 = \{1,2,3\}$};

  \draw[rounded corners=6pt,fill=targetgreen!6,draw=targetgreen!25,thick] (2.5,-0.2) rectangle (4.5,2.1);
  \node[font=\scriptsize\bfseries,text=targetgreen] at (3.5,2.35) {$C_2 = \{4,5,6\}$};

  \draw[rounded corners=6pt,fill=quotientorange!6,draw=quotientorange!25,thick] (5.5,-0.2) rectangle (7.5,2.1);
  \node[font=\scriptsize\bfseries,text=quotientorange] at (6.5,2.35) {$C_3 = \{7,8,9\}$};

  % Vertices
  \node[v1] (1) at (0,0.9)   {1};
  \node[v1] (2) at (0.9,1.6) {2};
  \node[v1] (3) at (0.9,0.2) {3};

  \node[v2] (4) at (3.0,0.9) {4};
  \node[v2] (5) at (3.9,1.6) {5};
  \node[v2] (6) at (3.9,0.2) {6};

  \node[v3] (7) at (6.0,0.9) {7};
  \node[v3] (8) at (6.9,1.6) {8};
  \node[v3] (9) at (6.9,0.2) {9};

  % Intra-triangle edges
  \path[intra] (1) edge (2) (1) edge (3) (2) edge (3);
  \path[intra] (4) edge (5) (4) edge (6) (5) edge (6);
  \path[intra] (7) edge (8) (7) edge (9) (8) edge (9);

  % Bridges
  \path[bridge] (3) edge node[above=1pt,font=\tiny\bfseries,text=trapred] {$(3,4)$} (4);
  \path[bridge] (6) edge node[above=1pt,font=\tiny\bfseries,text=trapred] {$(6,7)$} (7);
\end{tikzpicture}
\caption{The three-triangle chain: three internally dense clusters $C_1=\{1,2,3\}$, $C_2=\{4,5,6\}$, and
$C_3=\{7,8,9\}$ linked only by sparse bottleneck bridge edges $(3,4)$ and $(6,7)$ (red).}
\label{ch33:fig-chain}
\end{figure}

\begin{figure}[H]
\centering
\begin{tikzpicture}[scale=1.1,
   pt1/.style={circle,fill=sourceblue,draw=white,line width=0.8pt,inner sep=0pt,minimum size=5.5pt},
   pt2/.style={circle,fill=targetgreen,draw=white,line width=0.8pt,inner sep=0pt,minimum size=5.5pt},
   pt3/.style={circle,fill=quotientorange,draw=white,line width=0.8pt,inner sep=0pt,minimum size=5.5pt}]

  % Coordinate axes
  \draw[->,edgegray,thick] (-3.2,0) -- (3.2,0) node[right,font=\footnotesize] {$u_2$ (Fiedler)};
  \draw[->,edgegray,thick] (0,-2.3) -- (0,2.6) node[above,font=\footnotesize] {$u_3$};

  % Unit circle (Ng-Jordan-Weiss spherical normalization)
  \draw[edgegray!40,dashed,thick] (0,0) circle (2.2);
  \node[font=\tiny,text=edgegray,above right] at (1.55,1.55) {unit sphere $S^1$};

  % Cluster 1: negative u2, negative u3
  \draw[sourceblue!40,fill=sourceblue!10,rounded corners=8pt,thick] (-2.3,-1.4) rectangle (-1.2,-0.5);
  \node[pt1] at (-1.9,-0.8) {}; \node[font=\tiny,below left] at (-1.9,-0.8) {1};
  \node[pt1] at (-1.7,-1.1) {}; \node[font=\tiny,below] at (-1.7,-1.1) {2};
  \node[pt1] at (-1.4,-0.7) {}; \node[font=\tiny,above right] at (-1.4,-0.7) {3};
  \node[font=\scriptsize\bfseries,text=sourceblue,left=2pt] at (-2.3,-0.95) {$C_1=\{1,2,3\}$};

  % Cluster 2: near zero u2, positive u3
  \draw[targetgreen!40,fill=targetgreen!10,rounded corners=8pt,thick] (-0.6,1.4) rectangle (0.6,2.3);
  \node[pt2] at (-0.3,1.7) {}; \node[font=\tiny,left] at (-0.3,1.7) {4};
  \node[pt2] at (0.0,2.0)  {}; \node[font=\tiny,above] at (0.0,2.0) {5};
  \node[pt2] at (0.3,1.7)  {}; \node[font=\tiny,right] at (0.3,1.7) {6};
  \node[font=\scriptsize\bfseries,text=targetgreen,above=2pt] at (0,2.3) {$C_2=\{4,5,6\}$};

  % Cluster 3: positive u2, negative u3
  \draw[quotientorange!40,fill=quotientorange!10,rounded corners=8pt,thick] (1.2,-1.4) rectangle (2.3,-0.5);
  \node[pt3] at (1.4,-0.7) {}; \node[font=\tiny,above left] at (1.4,-0.7) {7};
  \node[pt3] at (1.7,-1.1) {}; \node[font=\tiny,below] at (1.7,-1.1) {8};
  \node[pt3] at (1.9,-0.8) {}; \node[font=\tiny,below right] at (1.9,-0.8) {9};
  \node[font=\scriptsize\bfseries,text=quotientorange,right=2pt] at (2.3,-0.95) {$C_3=\{7,8,9\}$};

  % Projection rays to the circle
  \draw[->,sourceblue!70,densely dotted,thick] (0,0) -- (-1.82,-1.23);
  \draw[->,targetgreen!70,densely dotted,thick] (0,0) -- (0.0,2.2);
  \draw[->,quotientorange!70,densely dotted,thick] (0,0) -- (1.82,-1.23);

  % Normalized points on circle
  \node[circle,draw=sourceblue,fill=sourceblue,inner sep=0pt,minimum size=4pt] at (-1.82,-1.23) {};
  \node[circle,draw=targetgreen,fill=targetgreen,inner sep=0pt,minimum size=4pt] at (0.0,2.2) {};
  \node[circle,draw=quotientorange,fill=quotientorange,inner sep=0pt,minimum size=4pt] at (1.82,-1.23) {};

  % Informational note
  \node[font=\scriptsize\itshape,text=edgegray!90] at (0,-2.0)
    {Radial normalization $y_i = U_{i,\cdot}/\|U_{i,\cdot}\|$ projects clusters onto isolated angular poles};
\end{tikzpicture}
\caption{2D spectral embedding of the three-triangle chain onto eigenvectors $u_2$ and $u_3$ of the
normalized Laplacian. The nine vertices collapse into three tight, well-separated point clouds.
Row-normalization projects each cluster onto the unit circle (dotted rays), where ordinary $k$-means
cleanly isolates the three communities without ambiguity.}
\label{ch33:fig-blobs}
\end{figure}

\section{Complexity}
\label{ch33:complexity}

The cost is dominated by the eigenvector computation. For a graph with $m$ edges and target $k$ clusters,
an iterative bottom-eigenvector solver (Lanczos / LOBPCG) costs $\Oh(m)$ per matrix--vector product and
needs a number of iterations that grows with the conditioning, which depends on the spectral gap.
Building the Laplacian is $\Oh(m)$. The sweep in Algorithm~\ref{ch33:alg-bisection} sorts ($\Oh(n\log n)$)
and then updates $\cut$ and $\vol$ incrementally as each vertex moves across, so all $n-1$ sweep cuts cost
$\Oh(m)$ total, not $\Oh(nm)$. The $k$-means stage on $n$ points in $\R^k$ costs $\Oh(nk\cdot I)$ per run
for $I$ iterations.

\begin{complexitybox}
\textbf{Spectral bisection (Alg.~\ref{ch33:alg-bisection}).} Eigenpair $(\lambda_2,v_2)$ by Lanczos:
$\widetilde\Oh(m / \sqrt{\lambda_2})$ to constant accuracy. Sweep rounding: $\Oh(m + n\log n)$.
\\[2pt]
\textbf{$k$-way (Alg.~\ref{ch33:alg-kway}).} Bottom $k$ eigenvectors: $\widetilde\Oh(mk/\sqrt{\mathrm{gap}})$.
Embedding and row-normalize: $\Oh(nk)$. $k$-means: $\Oh(nkI)$ per restart.
\\[2pt]
\textbf{Exact eigendecomposition} (dense, small graphs): $\Oh(n^3)$ -- what we used by hand on $\Gph_0$.
\\[2pt]
With nearly-linear-time Laplacian solvers (Spielman--Teng and successors), the eigenvector step can be made
$\widetilde\Oh(m)$, so the whole pipeline is near-linear in the graph size.
\end{complexitybox}

\section{Historical Context, Controversies, and the Road Ahead}
\label{ch33:history-context}

\begin{historybox}
In the early 1970s, two communities working on opposite problems hit the same mathematical wall. In Prague, Miroslav Fiedler was studying the physical mechanics of elastic trusses and electrical networks. He discovered that the second-smallest eigenvalue $\lambda_2$ of the graph Laplacian vanished if and only if the network fell into disconnected pieces, and its size measured structural cohesion. The signs of this eigenvector (the Fiedler vector) split the vertices into two natural clusters. At the exact same time at IBM, Donath and Hoffman were trying to partition logic circuits across silicon chips to minimize inter-chip wire delays. Because balanced partitioning is NP-hard, they turned to the eigenvalues of the adjacency and Laplacian matrices for lower bounds.

Three decades later, in 2000, computer vision researchers encountered this same geometry when trying to segment digital photographs into distinct objects. Local edge-detection filters routinely failed because pixel noise broke contours apart. Jianbo Shi and Jitendra Malik realized an image is better viewed as a global similarity graph, where cut boundaries should minimize the normalized cut. Soon after, Andrew Ng, Michael Jordan, and Yair Weiss developed the modern normalized spectral clustering pipeline, proving how eigenvector embeddings reveal geometric structure that standard methods miss.

Early attempts had stumbled into two seductive traps. The first was reaching for the polynomial-time min-cut algorithm from network flows, which in practice routinely isolates a single disconnected outlier or leaf vertex rather than splitting the graph into two meaningful halves. Penalizing cut sizes to avoid this imbalance instantly makes the problem NP-hard. The second trap was running classic $k$-means directly on raw spatial coordinates, which assumes clusters are spherical and convex, failing completely on intertwined spirals or concentric rings. Spectral clustering bypasses both traps: continuous eigenvector relaxation gives a smooth geometric embedding where simple thresholding recovers deep non-convex partitions, backed by Cheeger's discrete isoperimetric inequalities.
\end{historybox}

\subsection{What We Can Now Do, Honest Limitations, and the Road Ahead}

\begin{focusbox}
\textbf{Capabilities Unlocked}:
\begin{itemize}
  \item \textbf{Non-Convex Manifold Clustering}: By mapping vertices into the eigenspace of $\Lap_{\mathrm{sym}}$, complex non-linear
        structures (like concentric rings) are untangled into radially separated rays on the unit sphere, where simple $k$-means succeeds.
  \item \textbf{Global Relaxation}: Rather than getting trapped in local minima like greedy graph search, spectral methods solve for
        the global optimum of a continuous quadratic relaxation in polynomial time.
\end{itemize}

\textbf{Honest Limitations and Downsides}:
\begin{itemize}
  \item \textbf{Computational Burden on Large Graphs}: Full exact eigendecomposition requires $\Oh(n^3)$ operations. While sparse
        iterative solvers (Lanczos, LOBPCG) compute the bottom $k$ eigenvectors in $\widetilde\Oh(mk/\sqrt{\mathrm{gap}})$, they
        stall when the eigengap is tiny.
  \item \textbf{The Cheeger Quadratic Loss}: When $\lambda_2$ is small, the Cheeger guarantee $\phi \le \sqrt{2\lambda_2}$ can be loose.
        A graph with $\lambda_2 = 0.0001$ has conductance bounded only by $\sqrt{0.0002} \approx 0.014$, leaving room for an order-of-magnitude
        discrepancy between the continuous relaxation and the discrete cut.
  \item \textbf{Fragility to Degree Outliers}: In unnormalized clustering ($\Lap$), a single high-degree hub can dominate the quadratic
        form and pull the Fiedler vector toward itself, isolating the hub rather than uncovering balanced communities.
\end{itemize}

\textbf{Where We Go Next in the Book}:
\begin{itemize}
  \item In \textbf{Chapter~34 (Expander Graphs)}, we study the opposite extreme: graphs where $\lambda_2$ is bounded strictly away
        from zero, guaranteeing that \emph{no} sparse cut exists and information diffuses at maximum speed.
  \item In \textbf{Chapter~35 (Random Walks and Electrical Networks)}, we uncover the physical engine driving spectral clustering:
        the random-walk transition matrix $\mathbf{P} = \mathbf{D}^{-1}\Adjm$ and its Dirichlet potential theory.
  \item In \textbf{Chapter~41 (Laplacian Solvers and Electrical Flows)}, we learn how Spielman and Teng bypassed the $\Oh(n^3)$ barrier
        entirely, solving Laplacian linear systems in nearly-linear $\widetilde\Oh(m)$ time.
\end{itemize}
\end{focusbox}

\begin{applicationbox}
Spectral clustering shows up wherever data forms a similarity graph. \textbf{Image segmentation}: pixels are
vertices, edge weights are color/proximity similarities, NCut carves out objects (Shi--Malik).
\textbf{Community detection} in social and biological networks. \textbf{Mesh and circuit partitioning} for
parallel computation and VLSI layout, where a balanced cut minimizes inter-processor communication.
\textbf{Manifold learning / dimensionality reduction}: Laplacian eigenmaps are literally the spectral
embedding of Algorithm~\ref{ch33:alg-kway} used as features. \textbf{Document and gene clustering} via
$k$-nearest-neighbor similarity graphs.
\end{applicationbox}

\section{Building the graph: from data to similarity}
\label{ch33:graph-construction}

Often the input is not a graph but $n$ data points $x_1,\dots,x_n$. The clustering quality then hinges on
how you build the similarity graph -- a step at least as important as the spectral math.

\begin{itemize}
  \item \textbf{$\eps$-neighborhood graph}: connect $i,j$ if $\norm{x_i-x_j}\le\eps$. Unweighted, symmetric;
        sensitive to $\eps$ and to varying density.
  \item \textbf{$k$-nearest-neighbor graph}: connect $i$ to its $k$ nearest points (then symmetrize by OR or
        AND). Adapts to local density; the usual default.
  \item \textbf{Fully connected Gaussian (RBF) graph}: $w_{ij}=\exp(-\norm{x_i-x_j}^2/2\sigma^2)$. Smooth;
        the bandwidth $\sigma$ sets the scale of ``similar.''
\end{itemize}

\begin{trapbox}
The bandwidth $\sigma$ and the neighbor count $k$ are the real hyperparameters of spectral clustering, more
consequential than anything in the eigensolver. Too large and everything connects to everything ($\lambda_2$
big, no clusters visible); too small and the graph fragments into singletons ($\lambda_2\to0$, spurious
clusters). A robust default is to set $\sigma$ to the median pairwise distance among $k$-nearest neighbors,
or to use a self-tuning per-point bandwidth (Zelnik-Manor--Perona). When clusters fail to appear, suspect
the graph construction before the eigenvectors.
\end{trapbox}

\begin{connectionbox}
The whole chapter is one long bridge to the linear algebra companion. The relaxations are
\emph{Rayleigh-quotient minimizations}; the guarantees are \emph{Courant--Fischer} and \emph{Davis--Kahan};
the operator is the \emph{graph Laplacian} (Chapter~31) whose spectrum (Chapter~32) we read here. Forward, it
feeds \emph{expander graphs} (Chapter~34, where small conductance is the enemy and we want $\lambda_2$
\emph{large}), \emph{random walks and electrical networks} (Chapter~35, where $\Lap_{\mathrm{rw}}$ is the
walk generator and $\lambda_2^{\mathrm{nrm}}$ is the mixing rate), and \emph{PageRank} (Chapter~36, a
personalized random walk that is itself a localized spectral clustering method). To the probability/ML book:
$k$-means, Gaussian similarity kernels, and Laplacian eigenmaps for manifold learning.
\end{connectionbox}

\section{Exercises}
\label{ch33:exercises}

\begin{exercise}[Verify the Fiedler trace]
Compute the full eigendecomposition of the Laplacian of $\Gph_0$ (Figure~\ref{ch33:fig-running}) by hand or
by computer. Confirm $\lambda_2\approx0.44$ and that the Fiedler vector is antisymmetric under the
$L\leftrightarrow R$ mirror. Which eigenvalue and eigenvector correspond to the \emph{symmetric} mode, and
what cut, if any, does that mode suggest?
\end{exercise}

\begin{exercise}[The barbell and the dumbbell]
Let $\Gph$ be two cliques $K_t$ joined by a single edge (a barbell). Show $\lambda_2 = \Th(1/t^2)$ and that
the Fiedler vector cleanly separates the two cliques. Then replace the joining edge by a path of length
$\ell$ (a dumbbell). What happens to $\lambda_2$ and to the Fiedler partition as $\ell$ grows? Where does the
sweep cut the path?
\end{exercise}

\begin{exercise}[Prove the easy direction of Cheeger]
Show $\phi(\Gph)\ge\lambda_2^{\mathrm{nrm}}/2$ directly: plug the indicator-derived test vector $h$ of the
optimal conductance set into the Rayleigh quotient of $\Lap_{\mathrm{sym}}$ and bound it. Identify exactly
where the factor of $2$ is lost.
\end{exercise}

\begin{exercise}[RatioCut vs.\ NCut diverge]
Construct a graph on which the optimal RatioCut bipartition and the optimal NCut bipartition are
\emph{different}. (Hint: make one side have many low-degree vertices and the other few high-degree vertices,
so cardinality balance and volume balance disagree.) Compute both objectives for both partitions.
\end{exercise}

\begin{exercise}[Sweep beats sign]
Exhibit a connected graph whose Fiedler vector has a sign split that is a \emph{bad} cut, while some
non-sign sweep cut is good. Conclude that the sign threshold of Algorithm~\ref{ch33:alg-bisection} is not
safe to use alone, justifying the full sweep.
\end{exercise}

\begin{exercise}[Disconnected $=$ degenerate]
Show that if $\Gph$ has exactly $c$ connected components, then $\lambda_1=\dots=\lambda_c=0$ and the
$0$-eigenspace is spanned by component indicators. Explain why $k$-way spectral clustering with $k=c$ then
recovers the components \emph{exactly}, for any choice of orthonormal basis in line 2 of
Algorithm~\ref{ch33:alg-kway}.
\end{exercise}

\begin{exercise}[Implement and break]
Implement Algorithm~\ref{ch33:alg-kway}. Run it on three Gaussian blobs (easy), then on two concentric
circles (the classic case where $k$-means alone fails but spectral clustering with a $k$-NN graph
succeeds). Now \emph{break} it: shrink $\sigma$ until the inner circle fragments. Report the $\lambda$
spectrum at the breaking point and explain it via Section~\ref{ch33:graph-construction}.
\end{exercise}

\begin{exercise}[Eigengap heuristic]
For the three-triangle chain of Section~\ref{ch33:kway}, compute the bottom five eigenvalues and exhibit the
gap between $\lambda_3$ and $\lambda_4$. Then add edges between the triangles until the gap closes; relate
the closing gap to the merging of clusters and to Cheeger's inequality.
\end{exercise}

\begin{exercise}[Recursive vs.\ simultaneous]
Compare recursive spectral bisection (bisect, then bisect each half) against the simultaneous $k$-eigenvector
Algorithm~\ref{ch33:alg-kway} on a graph with three clusters of \emph{very} different sizes. Find an instance
where recursion makes the wrong first cut and never recovers, but the simultaneous method succeeds.
\end{exercise}

\begin{exercise}[Normalized vs.\ unnormalized perturbation]
Take a graph with a few very high-degree ``hub'' vertices. Show empirically that unnormalized spectral
clustering ($\Lap$) tends to isolate hubs while normalized clustering ($\Lap_{\mathrm{sym}}$) does not.
Explain using the role of $\mathbf D$ in the relaxation of Section~\ref{ch33:ncut}.
\end{exercise}
